\documentclass[11pt,a4paper]{article}
\usepackage{fontspec}\usepackage[french]{babel}
\usepackage[margin=1.8cm]{geometry}\usepackage{tabularx,colortbl,xcolor,array}
\setmainfont{PlusJakartaSans-Medium.ttf}[Path=fonts/,BoldFont=PlusJakartaSans-Bold.ttf]
\usepackage{xeCJK}\setCJKmainfont{WenQuanYiZenHei-subset.ttf}[Path=fonts/]
\pagestyle{empty}\setlength{\parindent}{0pt}
\begin{document}
{\LARGE\bfseries Chinois — Première A\par}\smallskip
{\large Répartition des 43 leçons par chapitre et par trimestre\par}\medskip
Cours « signés OnBuch+ », conformes à la fiche de progression harmonisée du MINESEC (année scolaire 2026-2027). Chaque leçon est un PDF autonome (3736 pages au total pour les 43 leçons).\par\medskip
\textbf{Contenu type d'une leçon :} page de garde, sommaire, objectifs, prérequis, sections de cours (textes en chinois simplifié avec pinyin et traduction, vocabulaire en tableaux, structures de grammaire, clés et ordre des traits, exemples traduits), activité d'intégration, méthodes et exercices résolus, exercices (connaissances, entraînement, préparation à l'examen), corrigés, fiche bilan.\par\medskip
\textbf{Les 5 modules :} Module 1 : Vie familiale et intégration sociale (leçons 1--8) ; Module 2 : Environnement et santé (leçons 9--18) ; Module 3 : Citoyenneté et ouverture au monde (leçons 19--30) ; Module 4 : Activités économiques et monde du travail (leçons 31--34) ; Module 5 : Mass médias et télécommunication (leçons 35--43).\par\bigskip
\section*{Module 1 : Vie familiale et intégration sociale}
\textit{8 leçons, leçons 1 à 8, 745 pages.}\par\medskip
\begin{tabularx}{\linewidth}{|c|X|l|l|c|}\hline
\rowcolor{gray!20} N° & Titre de la leçon & Trimestre & Période & Pages \\ \hline
1 & Leçon 1 — Vocabulaire et compréhension de texte : amour de la famille et la solidarité familiale \newline {\small\emph{Type : vocabulaire — fichier : ZHA01-vocabulaire-amour-de-la-famille-et-la-solidarite-f.pdf}} & 1er trimestre & 07/09 – 23/10/2026 & 77 \\ \hline
2 & Leçon 2 — Grammaire : 第一……第二…. ; 不但……而且…. \newline {\small\emph{Type : grammaire — fichier : ZHA02-grammaire.pdf}} & 1er trimestre & 07/09 – 23/10/2026 & 80 \\ \hline
3 & Leçon 3 — Expression orale : amour de la famille et la solidarité familiale \newline {\small\emph{Type : oral — fichier : ZHA03-oral-amour-de-la-famille-et-la-solidarite-familial.pdf}} & 1er trimestre & 07/09 – 23/10/2026 & 109 \\ \hline
4 & Leçon 4 — Expression écrite : écriture des caractères de l'amour et de la famille ; types de phrases étudiées \newline {\small\emph{Type : ecrit — fichier : ZHA04-ecrit-ecriture-des-caracteres-de-l-amour-et-de-la.pdf}} & 1er trimestre & 07/09 – 23/10/2026 & 96 \\ \hline
5 & Leçon 5 — Vocabulaire et compréhension de texte : délinquance juvénile \newline {\small\emph{Type : vocabulaire — fichier : ZHA05-vocabulaire-delinquance-juvenile.pdf}} & 1er trimestre & 07/09 – 23/10/2026 & 86 \\ \hline
6 & Leçon 6 — Grammaire : 如果……，那么……；les nombres ；有的……有的……；是……的 \newline {\small\emph{Type : grammaire — fichier : ZHA06-grammaire-les-nombres-les-nombres.pdf}} & 1er trimestre & 07/09 – 23/10/2026 & 87 \\ \hline
7 & Leçon 7 — Expression orale : délinquance juvénile \newline {\small\emph{Type : oral — fichier : ZHA07-oral-delinquance-juvenile.pdf}} & 1er trimestre & 07/09 – 23/10/2026 & 116 \\ \hline
8 & Leçon 8 — Expression écrite : la délinquance juvénile \newline {\small\emph{Type : ecrit — fichier : ZHA08-ecrit-la-delinquance-juvenile.pdf}} & 1er trimestre & 07/09 – 23/10/2026 & 94 \\ \hline
\end{tabularx}\bigskip
\section*{Module 2 : Environnement et santé}
\textit{10 leçons, leçons 9 à 18, 826 pages.}\par\medskip
\begin{tabularx}{\linewidth}{|c|X|l|l|c|}\hline
\rowcolor{gray!20} N° & Titre de la leçon & Trimestre & Période & Pages \\ \hline
9 & Leçon 9 — Vocabulaire et compréhension de texte : importance de la pratique du sport \newline {\small\emph{Type : vocabulaire — fichier : ZHA09-vocabulaire-importance-de-la-pratique-du-sport.pdf}} & 1er trimestre & 26/10 – 18/12/2026 & 71 \\ \hline
10 & Leçon 10 — Grammaire : Phrases comparatives avec…不如… ; A比B 形容词+一点ou 多了 ; A比B更/ 还+形容词 ; 为了 ; 不但…而且… \newline {\small\emph{Type : grammaire — fichier : ZHA10-grammaire-phrases-comparatives-avec-ab-ou-ab.pdf}} & 1er trimestre & 26/10 – 18/12/2026 & 72 \\ \hline
11 & Leçon 11 — Expression orale : importance de la pratique du sport \newline {\small\emph{Type : oral — fichier : ZHA11-oral-importance-de-la-pratique-du-sport.pdf}} & 1er trimestre & 26/10 – 18/12/2026 & 117 \\ \hline
12 & Leçon 12 — Expression écrite : caractères liés à l'importance de la pratique du sport \newline {\small\emph{Type : ecrit — fichier : ZHA12-ecrit-caracteres-lies-a-l-importance-de-la-pratiqu.pdf}} & 1er trimestre & 26/10 – 18/12/2026 & 68 \\ \hline
13 & Leçon 13 — Vocabulaire et compréhension de texte : maladies en général et du SIDA en particulier \newline {\small\emph{Type : vocabulaire — fichier : ZHA13-vocabulaire-maladies-en-general-et-du-sida-en-part.pdf}} & 1er trimestre & 26/10 – 18/12/2026 & 89 \\ \hline
14 & Leçon 14 — Grammaire : Le compl. de résultat ; fraction avec 分之、百分之、三分之 \newline {\small\emph{Type : grammaire — fichier : ZHA14-grammaire-le-compl-de-resultat-fraction-avec.pdf}} & 1er trimestre & 26/10 – 18/12/2026 & 75 \\ \hline
15 & Leçon 15 — Vocabulaire et compréhension de texte : protection de l’environnement \newline {\small\emph{Type : vocabulaire — fichier : ZHA15-vocabulaire-protection-de-lenvironnement.pdf}} & 1er trimestre & 26/10 – 18/12/2026 & 75 \\ \hline
16 & Leçon 16 — Grammaire : Les compl.de durée ; redoublement des adjectifs \newline {\small\emph{Type : grammaire — fichier : ZHA16-grammaire-les-compl-de-duree-redoublement-des-adje.pdf}} & 1er trimestre & 26/10 – 18/12/2026 & 82 \\ \hline
17 & Leçon 17 — Expression orale : protection de l’environnement \newline {\small\emph{Type : oral — fichier : ZHA17-oral-protection-de-lenvironnement.pdf}} & 1er trimestre & 26/10 – 18/12/2026 & 102 \\ \hline
18 & Leçon 18 — Expression écrite : protection de l'environnement (你应该怎么保护学校的环境) \newline {\small\emph{Type : ecrit — fichier : ZHA18-ecrit-protection-de-l-environnement.pdf}} & 1er trimestre & 26/10 – 18/12/2026 & 75 \\ \hline
\end{tabularx}\bigskip
\section*{Module 3 : Citoyenneté et ouverture au monde}
\textit{12 leçons, leçons 19 à 30, 1046 pages.}\par\medskip
\begin{tabularx}{\linewidth}{|c|X|l|l|c|}\hline
\rowcolor{gray!20} N° & Titre de la leçon & Trimestre & Période & Pages \\ \hline
19 & Leçon 19 — Vocabulaire et compréhension de texte : vie culturelle au Cameroun et en Chine \newline {\small\emph{Type : vocabulaire — fichier : ZHA19-vocabulaire-vie-culturelle-au-cameroun-et-en-chine.pdf}} & 2e trimestre & 04/01 – 19/03/2027 & 73 \\ \hline
20 & Leçon 20 — Grammaire : 比较字句:……跟/和…. (不)一样et ……跟/和…. (不)一样+ 形容词 ; 一方面… ,另一方面… ; 如果…, 但是/可是… \newline {\small\emph{Type : grammaire — fichier : ZHA20-grammaire-et.pdf}} & 2e trimestre & 04/01 – 19/03/2027 & 85 \\ \hline
21 & Leçon 21 — Expression orale : vie culturelle au Cameroun et en Chine \newline {\small\emph{Type : oral — fichier : ZHA21-oral-vie-culturelle-au-cameroun-et-en-chine.pdf}} & 2e trimestre & 04/01 – 19/03/2027 & 83 \\ \hline
22 & Leçon 22 — Expression écrite : culture chinoise et camerounaise (你喜欢喀麦隆还是中国艺术？为什么？) \newline {\small\emph{Type : ecrit — fichier : ZHA22-ecrit-culture-chinoise-et-camerounaise.pdf}} & 2e trimestre & 04/01 – 19/03/2027 & 100 \\ \hline
23 & Leçon 23 — Vocabulaire et compréhension de texte : exode rural \newline {\small\emph{Type : vocabulaire — fichier : ZHA23-vocabulaire-exode-rural.pdf}} & 2e trimestre & 04/01 – 19/03/2027 & 88 \\ \hline
24 & Leçon 24 — Grammaire : 越来越； ; 一方面……,另一方面…… \newline {\small\emph{Type : grammaire — fichier : ZHA24-grammaire.pdf}} & 2e trimestre & 04/01 – 19/03/2027 & 86 \\ \hline
25 & Leçon 25 — Expression orale : exode rural \newline {\small\emph{Type : oral — fichier : ZHA25-oral-exode-rural.pdf}} & 2e trimestre & 04/01 – 19/03/2027 & 92 \\ \hline
26 & Leçon 26 — Expression écrite : caractères complexes liés à l'exode rural \newline {\small\emph{Type : ecrit — fichier : ZHA26-ecrit-caracteres-complexes-lies-a-l-exode-rural.pdf}} & 2e trimestre & 04/01 – 19/03/2027 & 96 \\ \hline
27 & Leçon 27 — Vocabulaire et compréhension de texte : corruption \newline {\small\emph{Type : vocabulaire — fichier : ZHA27-vocabulaire-corruption.pdf}} & 2e trimestre & 08/03 – 19/03/2027 & 83 \\ \hline
28 & Leçon 28 — Grammaire : 把字句 \newline {\small\emph{Type : grammaire — fichier : ZHA28-grammaire.pdf}} & 2e trimestre & 08/03 – 19/03/2027 & 81 \\ \hline
29 & Leçon 29 — Expression orale : corruption \newline {\small\emph{Type : oral — fichier : ZHA29-oral-corruption.pdf}} & 2e trimestre & 08/03 – 19/03/2027 & 66 \\ \hline
30 & Leçon 30 — Expression écrite : facteurs de la corruption ; caractères complexes \newline {\small\emph{Type : ecrit — fichier : ZHA30-ecrit-facteurs-de-la-corruption-caracteres-complex.pdf}} & 2e trimestre & 08/03 – 19/03/2027 & 113 \\ \hline
\end{tabularx}\bigskip
\section*{Module 4 : Activités économiques et monde du travail}
\textit{4 leçons, leçons 31 à 34, 314 pages.}\par\medskip
\begin{tabularx}{\linewidth}{|c|X|l|l|c|}\hline
\rowcolor{gray!20} N° & Titre de la leçon & Trimestre & Période & Pages \\ \hline
31 & Leçon 31 — Vocabulaire et compréhension de texte : économie \newline {\small\emph{Type : vocabulaire — fichier : ZHA31-vocabulaire-economie.pdf}} & 3e trimestre & 22/03 – 23/04/2027 & 64 \\ \hline
32 & Leçon 32 — Grammaire : compl. d’état \newline {\small\emph{Type : grammaire — fichier : ZHA32-grammaire-compl-detat.pdf}} & 3e trimestre & 22/03 – 23/04/2027 & 75 \\ \hline
33 & Leçon 33 — Expression orale : économie \newline {\small\emph{Type : oral — fichier : ZHA33-oral-economie.pdf}} & 3e trimestre & 22/03 – 23/04/2027 & 75 \\ \hline
34 & Leçon 34 — Expression écrite : économie du Cameroun (为什么喀麦隆经济比泰发展？) \newline {\small\emph{Type : ecrit — fichier : ZHA34-ecrit-economie-du-cameroun.pdf}} & 3e trimestre & 22/03 – 23/04/2027 & 100 \\ \hline
\end{tabularx}\bigskip
\section*{Module 5 : Mass médias et télécommunication}
\textit{9 leçons, leçons 35 à 43, 805 pages.}\par\medskip
\begin{tabularx}{\linewidth}{|c|X|l|l|c|}\hline
\rowcolor{gray!20} N° & Titre de la leçon & Trimestre & Période & Pages \\ \hline
35 & Leçon 35 — Vocabulaire et compréhension de texte : internet \newline {\small\emph{Type : vocabulaire — fichier : ZHA35-vocabulaire-internet.pdf}} & 3e trimestre & 26/04 – 11/06/2027 & 98 \\ \hline
36 & Leçon 36 — Grammaire : antonyme 新 / 旧 ; verbe modal 可以 \newline {\small\emph{Type : grammaire — fichier : ZHA36-grammaire-antonyme-verbe-modal.pdf}} & 3e trimestre & 26/04 – 11/06/2027 & 91 \\ \hline
37 & Leçon 37 — Expression orale : internet \newline {\small\emph{Type : oral — fichier : ZHA37-oral-internet.pdf}} & 3e trimestre & 26/04 – 11/06/2027 & 71 \\ \hline
38 & Leçon 38 — Expression écrite : origine du caractère 电 et radicaux ; caractères des réseaux sociaux \newline {\small\emph{Type : ecrit — fichier : ZHA38-ecrit-origine-du-caractere-et-radicaux-caracteres.pdf}} & 3e trimestre & 26/04 – 11/06/2027 & 109 \\ \hline
39 & Leçon 39 — Vocabulaire et compréhension de texte : rédaction d’une lettre formelle \newline {\small\emph{Type : vocabulaire — fichier : ZHA39-vocabulaire-redaction-dune-lettre-formelle.pdf}} & 3e trimestre & 26/04 – 11/06/2027 & 85 \\ \hline
40 & Leçon 40 — Grammaire : 一边……一边…….. ; la ponctuation en chinois \newline {\small\emph{Type : grammaire — fichier : ZHA40-grammaire-la-ponctuation-en-chinois.pdf}} & 3e trimestre & 26/04 – 11/06/2027 & 99 \\ \hline
41 & Leçon 41 — Expression orale : rédaction d’une lettre formelle \newline {\small\emph{Type : oral — fichier : ZHA41-oral-redaction-dune-lettre-formelle.pdf}} & 3e trimestre & 26/04 – 11/06/2027 & 87 \\ \hline
42 & Leçon 42 — Expression écrite : caractères complexes liés à la gestion des réseaux sociaux \newline {\small\emph{Type : ecrit — fichier : ZHA42-ecrit-caracteres-complexes-lies-a-la-gestion-des-r.pdf}} & 3e trimestre & 26/04 – 11/06/2027 & 81 \\ \hline
43 & Leçon 43 — Révisions générales \newline {\small\emph{Type : revision — fichier : ZHA43-revision-revisions.pdf}} & 3e trimestre & 26/04 – 11/06/2027 & 84 \\ \hline
\end{tabularx}\bigskip
\end{document}
